\section{Experiments}
\label{sec:experiments}

We evaluate generational depth estimation on ten model configurations spanning six architecture families:
Pythia-1.4B, Pythia-6.9B, OLMo-1B, Qwen2.5-1.5B-Instruct (rewrite and continuation), Qwen2.5-1.5B (base), Qwen2.5-7B, LLaMA-3.1-8B, Mistral-7B, and Gemma-2-2B, all in continuation mode except where noted.
All experiments use the OBD-LightGBM classifier with the 15-dimensional feature set (Section~\ref{sec:method}), with 5{,}000 texts per depth level (generations 0--5).

%% ====================================================================
\subsection{$\Kempirical \leq 2$ Across Model Configurations (C4)}
\label{sec:kstar_boundary}

\begin{table*}[t]
\centering
\small
\caption{Pairwise $\auc$ for adjacent depth pairs with bootstrap 95\% CIs and permutation $p$-values for boundary rows.
$\Kempirical = 2$ requires 1v2 $\auc > \theta{=}0.60$ and 2v3 $\auc < \theta$; CIs quantify boundary uncertainty.
Extended results (all models) in Table~\ref{tab:kstar_ci_extended}.}
\label{tab:kstar_ci}
\begin{tabular}{lcccccccccc}
\toprule
 & \multicolumn{2}{c}{Qwen-rw} & \multicolumn{2}{c}{Pythia} & \multicolumn{2}{c}{OLMo} & \multicolumn{2}{c}{Qwen-cont} & \multicolumn{2}{c}{Qwen-7B\rlap{$^\ddagger$}} \\
\cmidrule(lr){2-3}\cmidrule(lr){4-5}\cmidrule(lr){6-7}\cmidrule(lr){8-9}\cmidrule(lr){10-11}
 & $\auc$ & 95\% CI & $\auc$ & 95\% CI & $\auc$ & 95\% CI & $\auc$ & 95\% CI & $\auc$ & 95\% CI \\
\midrule
0v1 & .9956 & & .9932 & & .9860 & & .9971 & & .9638 & [.961,\,.967] \\
1v2 & .6955 & [.686,\,.704]\rlap{$^*$} & .6539 & [.645,\,.663]\rlap{$^*$} & .5080 & [.496,\,.519] & .6112 & [.602,\,.621]\rlap{$^*$} & .6168 & [.608,\,.626]\rlap{$^*$} \\
2v3 & .5893 & [.580,\,.599] & .5514 & [.542,\,.563] & .4973 & [.486,\,.509] & .5609 & [.551,\,.571] & .5492 & [.539,\,.559] \\
3v4 & .5435 & & .5152 & & .5024 & & .5409 & & .5276 & \\
4v5 & .5340 & & .5136 & & .4942 & & .5141 & & .5085 & \\
\bottomrule
\multicolumn{11}{l}{\footnotesize $^*$Permutation $p < 0.001$ for all 1v2 comparisons shown.} \\
\multicolumn{11}{l}{\footnotesize $^\ddagger$Generator: Qwen2.5-7B (fp16); scorer: Qwen2.5-1.5B.} \\
\multicolumn{11}{l}{\footnotesize Full results for all 10 configurations (including Qwen-7B, Qwen-base, Gemma-2-2B, and Pythia-6.9B) are provided in Table~\ref{tab:kstar_ci_extended}.}
\end{tabular}
\end{table*}

\begin{figure}[t]
\centering
\includegraphics[width=\columnwidth]{figures/fig2_auc_decay.pdf}
\caption{Pairwise $\auc$ decay across depth pairs for model-mode configurations (solid) and decoding strategies (dashed); shaded bands show bootstrap 95\% confidence intervals. The horizontal dashed line marks $\theta = 0.60$.}
\label{fig:auc_decay}
\end{figure}

Table~\ref{tab:kstar_ci} and Figure~\ref{fig:auc_decay} show $\Kempirical \leq 2$ across all 10 model configurations: 6/10 configurations yield $\Kempirical{=}2$ (2/6 families unanimously: Qwen with 4 configs, LLaMA-8B; Pythia-1.4B provides additional cross-family evidence) (1v2 $\auc$ range: 0.6112--0.6955; all CI lower bounds $> \theta$; permutation $p < 0.001$), while OLMo-1B, Mistral-7B, Pythia-6.9B, and Gemma-2-2B yield $\Kempirical{=}1$ (1v2 $\auc$: 0.508, 0.594, 0.596, and 0.598, all below $\theta$; Section~\ref{sec:boundary_cases}).
LLaMA-3.1-8B confirms $\Kempirical{=}2$ (0v1 $\auc = 0.9225$, 1v2 $\auc = 0.6123\;[0.603,\,0.622]$; CI lower bound above $\theta$; Table~\ref{tab:kstar_ci_extended}).
All four decoding strategies likewise yield $\Kempirical{=}2$ (Table~\ref{tab:s2_decoding}).
The 1v2 spread due to generation mode ($\Delta = 8.4$\,pp within Qwen) exceeds the spread across $\Kempirical{=}2$ continuation models ($\Delta = 4.3$\,pp).

\paragraph{Scale and instruction-tuning invariance.}
$\Kempirical{=}2$ persists from 1.5B through 8B: Qwen-7B (1v2 $\auc = 0.6168$ $[0.608, 0.626]$) and LLaMA-3.1-8B (1v2 $\auc = 0.6123$) both exceed $\theta$, despite progressively reduced 0v1 signal (0.9638 and 0.9225 vs.\ $> 0.99$ at $\leq$1.5B).
A base-vs.-instruct ablation at 1.5B likewise yields $\Kempirical = 2$ (1v2 $\auc = 0.626$), confirming the boundary is not an instruction-tuning artifact.

\begin{figure}[t]
\centering
\includegraphics[width=\columnwidth]{figures/fig_kstar_theta.pdf}
\caption{$\Kempirical(\theta)$ sensitivity curve for six configurations. $\Kempirical$: maximum discriminable depth; $\theta$: AUC decision threshold. Five of six configurations yield $\Kempirical{=}2$ at $\theta{=}0.60$ (dashed line); OLMo yields $\Kempirical{=}1$. At $\theta{=}0.70$, all drop to $\Kempirical{=}1$.}
\label{fig:kstar_theta}
\end{figure}

Figure~\ref{fig:kstar_theta} shows $\Kempirical$ as a function of $\theta$: five of six configurations yield $\Kempirical = 2$ within $\theta \in [0.51, 0.61]$ (OLMo: $\Kempirical = 1$ throughout); at $\theta{=}0.55$, two Qwen configurations reach $\Kempirical = 3$; at $\theta{=}0.70$, all drop to $\Kempirical = 1$.
All four decoding strategies likewise yield $\Kempirical = 2$ at $\theta{=}0.60$ (1v2 $\auc \in [0.687, 0.704]$; Table~\ref{tab:s2_decoding}).

%% ====================================================================
\subsection{Cross-Model Transfer}
\label{sec:transfer}

\begin{figure}[t]
\centering
\includegraphics[width=0.85\columnwidth]{figures/fig3_transfer_heatmap.pdf}
\caption{Cross-model transfer heatmap. Each cell shows 6-class ordinal accuracy when trained on the row model and tested on the column model, divided by in-domain accuracy (retention rate). Pythia$\leftrightarrow$OLMo: 90--92\%; Qwen-rewrite$\to$continuation: 61--63\%.}
\label{fig:transfer_heatmap}
\end{figure}

Continuation-mode classifiers transfer well across generators (Figure~\ref{fig:transfer_heatmap}; full matrix in Table~\ref{tab:s1_transfer}): Pythia$\leftrightarrow$OLMo retains 90--92\%, Qwen-cont$\leftrightarrow$Pythia 84--85\%, Qwen-cont$\leftrightarrow$OLMo 80--85\% (0v1 cross-model $\auc > 0.96$ for all pairs).
Qwen-rewrite classifiers transfer poorly (61--63\%), confirming reliance on mode-specific artifacts rather than generalizable depth features.

\paragraph{Cross-scorer validation.}
Replacing the self-scorer with Pythia-1.4B or OLMo-1B preserves feature rankings ($\rho > 0.90$) but attenuates absolute discriminability at the boundary (Table~\ref{tab:crossscorer}).
Feature ranking independence is a strong result: the same features dominate under all scorers.
$\Kempirical$ agreement, however, is signal-strength dependent: Qwen-rewrite (strong signal, 1v2 $\auc = 0.696$) retains $\Kempirical{=}2$ under all scorers, but Qwen-continuation (marginal, 1v2 $\auc = 0.611$, only 0.011 above $\theta$) drops to $\Kempirical{=}1$ (Pythia: $0.5923$; OLMo: $0.5858$).
A three-way ANOVA (Generator $\times$ Scorer $\times$ Depth-Pair; Table~\ref{tab:anova_crossscorer}) confirms a statistically significant but modest Scorer main effect ($\eta^2_p{=}.166$, $p < .001$) and a significant Generator$\times$Scorer interaction ($\eta^2_p{=}.146$), indicating that absolute AUC values are scorer-dependent even when $\Kempirical$ classifications largely agree.

%% ====================================================================
\subsection{Feature Ordering and Generation Mode}
\label{sec:feature_degradation}
\label{sec:generation_mode}

Among continuation-mode models, single-feature $\auc$ rankings are highly correlated ($\rho \geq 0.85$; Table~\ref{tab:top_features}).
The top-3 features (\texttt{surp\_d1\_std}, \texttt{surp\_std}, \texttt{surp\_d2\_std}) are shared across all three families, revealing that autoregressive generation progressively damps surprisal variability.
Qwen-rewrite diverges (\texttt{freq\_entropy} dominates), consistent with its poor cross-model transfer (61--63\%).
Within Qwen, rewrite yields 1v2 $\auc = 0.6955$ vs.\ continuation's $0.6112$ ($\Delta = +8.4$\,pp); continuation mode is therefore recommended as the canonical condition.

%% ====================================================================
\subsection{Neural Probe}
\label{sec:neural_probe}

Replacing the 15D features with frozen encoder hidden states (MLP probe; Table~\ref{tab:neural_probe_crossmodel}), Qwen achieves $\Kempirical^{\text{neural}} = 3$ (2v3 $\auc = 0.6159$), while Pythia remains at $\Kempirical^{\text{neural}} = 2$ ($\auc_{2v3} = 0.5963$); OLMo yields $\Kempirical^{\text{neural}} = 1$ (1v2 $\auc = 0.5054$, consistent with its 15D boundary).
The Qwen-only elevation reflects self-familiarity artifacts that transfer poorly (Section~\ref{sec:transfer}); the 15D features yield $\Kempirical = 2$ for most families (OLMo: $\Kempirical = 1$).

\begin{minipage}{\columnwidth}
\centering
\footnotesize
\captionof{table}{Neural probe pairwise AUC using frozen-encoder MLP. Each cell reports AUC for classifying depth $d$ vs.\ $d{+}1$ text. $\Kempirical^{\text{neural}}$: maximum depth where AUC exceeds $\theta = 0.60$.}
\label{tab:neural_probe_crossmodel}
\begin{tabular}{lccc}
\toprule
\textbf{Model} & \textbf{2v3 $\auc$} & \textbf{95\% CI} & $\Kempirical^{\text{neural}}$ \\
\midrule
Qwen-1.5B  & 0.6159 & $[0.610,\,0.622]$ & 3 \\
Pythia-1.4B & 0.5963 & $[0.589,\,0.603]$ & 2 \\
OLMo-1B    & 0.5056 & $[0.497,\,0.516]$ & 1\rlap{$^\dagger$} \\
\bottomrule
\multicolumn{4}{l}{\footnotesize $^\dagger$OLMo 1v2 $\auc = 0.5054$ (below $\theta$); $\Kempirical^{\text{neural}} = 1$.}
\end{tabular}
\end{minipage}

%% ====================================================================
\subsection{Cross-Domain Variation}
\label{sec:cross_domain}

The preceding experiments use C4 (general web text).
To test domain dependence, we evaluate on Wikipedia ($n{=}5{,}000$) and arXiv abstracts ($n{=}2{,}200$) using two model families.
Wikipedia yields $\Kempirical{=}2$ (1v2 $\auc = 0.623$, 2v3 $\auc = 0.571$), matching the C4 baseline.
Both Qwen-1.5B and Pythia-1.4B yield $\Kempirical{=}3$ for arXiv (Table~\ref{tab:domain_kstar}), demonstrating that structured academic text permits discrimination to greater depth independently of architecture.
The elevated signal reflects higher per-generation distributional shift (1v2 JSD $= 0.077$ vs.\ ${\sim}0.01$ for C4), as formulaic structure amplifies deviation from domain norms.

\begin{minipage}{\columnwidth}
\centering
\footnotesize
\captionof{table}{Cross-domain $\Kempirical$ comparison. Rows: generator model and text domain (C4, Wikipedia, arXiv). $\Kempirical{=}3$ for arXiv indicates structured text preserves discriminability to greater depth.}
\label{tab:domain_kstar}
\begin{tabular}{llccc}
\toprule
\textbf{Domain} & \textbf{Model} & \textbf{1v2 $\auc$} & \textbf{2v3 $\auc$} & $\Kempirical$ \\
\midrule
arXiv & Qwen-1.5B  & 0.858 & 0.6436 & 3 \\
arXiv & Pythia-1.4B & 0.805 & 0.6545 & 3 \\
C4    & Qwen-1.5B  & 0.63  & 0.56   & 2 \\
Wikipedia & Qwen-1.5B & 0.62 & 0.55 & 2 \\
\bottomrule
\end{tabular}
\end{minipage}

%% ====================================================================
\subsection{Binary Detector Confidence}
\label{sec:binary_confidence}

Binary detectors (DivEye~\cite{basani2026diveye}, Binoculars~\cite{hans2024binoculars}) show moderate rank correlations with depth ($|\rho| = 0.06$--$0.46$; Table~\ref{tab:s9_binary_spearman}) but fail as depth proxies: 6-class accuracy is at or below chance ($\leq 0.199$ vs.\ 0.167 baseline), pairwise $\auc$ yields $\Kempirical^{\text{binary}} \leq 1$, and multiple conditions exhibit non-monotonic AUC at deeper pairs.
The observed correlation is an artifact of 0v1 separation rather than fine-grained depth signal.
